\chapter{Теоретические сведения}

Данная глава вводит понятия, которые дальше используются в методике и в интерпретации результатов. Здесь рассматриваются внутренние представления нейронных сетей, частотное описание изображений и карт признаков, радиальные спектры, круг Найквиста, а также архитектуры ResNet, Vision Transformer и Swin Transformer. Отдельно оговариваются ограничения спектральной интерпретации: спектр показывает распределение энергии по пространственным частотам, но сам по себе не доказывает, какие признаки являются полезными для классификации.

\section{Внутренние представления и частотная интерпретация}

\subsection{Внутренние представления моделей компьютерного зрения}

Глубокую нейронную сеть можно рассматривать как последовательность преобразований входного изображения. На вход подаётся тензор пикселей, на выходе получается предсказание класса или другой целевой результат. Между ними находятся промежуточные слои. Их выходы обычно называют внутренними представлениями, признаками или активациями.

Внутреннее представление показывает, в какой форме модель хранит информацию об изображении на некотором этапе обработки. В ранних слоях признаки часто ближе к исходному изображению: они могут реагировать на локальные границы, контуры, текстуры и цветовые переходы. В более глубоких слоях признаки обычно становятся менее привязанными к отдельным пикселям и могут описывать более крупные части объектов или повторяющиеся структуры. Это не означает, что каждый глубокий слой всегда содержит только «семантику», а каждый ранний — только «текстуру». Такая граница грубая. Но как рабочее описание она полезна.

Для свёрточных сетей внутренние представления имеют вид карт признаков. Карта признаков — это двумерная сетка значений, заданная для каждого канала. Если исходное изображение содержит три цветовых канала, то глубокий слой сети может содержать сотни или тысячи каналов. Каждый канал можно рассматривать как отдельный тип отклика, который сеть сформировала в процессе обучения.

В трансформерных архитектурах представления устроены иначе. Изображение разбивается на патчи, после чего каждый патч переводится в вектор — токен. Дальше модель работает не с плотной картой пикселей, а с последовательностью токенов. Чтобы применить к таким представлениям двумерный спектральный анализ, токены возвращают в решётку, соответствующую расположению патчей в исходном изображении. Это удобно, но здесь есть нюанс: такая решётка не является полной копией карты признаков CNN. Один токен уже содержит агрегированную информацию о целом фрагменте изображения.

Классическая работа \cite{yosinski2014transferable} показывает, что признаки разных уровней обладают разной переносимостью между задачами. Ранние признаки часто более универсальны, поздние сильнее зависят от конкретного набора данных и задачи. В данной работе этот взгляд дополняется частотной интерпретацией: нас интересует, как по глубине модели меняется распределение энергии внутренних представлений по пространственным частотам.

\subsection{Изображение как пространственный сигнал}

Изображение можно рассматривать как пространственный сигнал, заданный на двумерной координатной сетке. В каждой точке этой сетки находится значение яркости или цвета. Если при движении по изображению значения меняются плавно, это соответствует низким пространственным частотам. Если значения резко меняются от точки к точке, речь идёт о более высоких частотах.

Плавный фон, крупный силуэт объекта и большие однородные области обычно дают вклад в низкие частоты. Резкие границы, тонкие линии, мелкие текстуры и шум связаны с более высокими частотами. Реальное изображение почти всегда содержит оба типа компонент. Например, фотография животного содержит низкочастотную информацию о форме тела и фоне, а также более высокочастотную информацию о шерсти, краях и мелких деталях.

В задачах компьютерного зрения это различие имеет смысл. Для грубого различения объектов часто достаточно общей формы. Для задач, где решают мелкие отличия, высокочастотные компоненты могут быть существенными. При этом нельзя автоматически считать высокие частоты «полезными», а низкие — «недостаточными». В высоких частотах может находиться и шум, а низкие частоты могут давать устойчивое описание крупной структуры. Поэтому дальше частотный анализ используется как способ описания представлений, а не как готовая оценка их качества.

\subsection{Дискретное преобразование Фурье, DC-компонента и спектр мощности}

Спектр описывает сигнал через набор частотных компонент. В пространственной области мы смотрим, как значения расположены по координатам. В частотной области — из каких изменений разного масштаба состоит сигнал. Для двумерных изображений и карт признаков используется двумерное дискретное преобразование Фурье.

В спектральном анализе различают амплитуду, фазу и энергию. Амплитуда показывает силу частотной компоненты, фаза связана с её пространственным положением, а энергия описывает вклад частоты в общий сигнал. В данной работе анализируется спектр мощности, то есть квадрат модуля коэффициентов Фурье. Такой выбор связан с целью исследования: нужно сравнить не точное положение структур в пространстве, а распределение энергии между частотными диапазонами.

Отдельно нужно упомянуть нулевую частоту, или DC-компоненту. Она соответствует постоянной составляющей сигнала, то есть его среднему значению. В картах признаков DC-компонента может быть очень большой и визуально скрывать различия между ненулевыми частотами. Поэтому перед преобразованием Фурье из каждой карты признаков вычитается пространственное среднее. В методической главе этот шаг описан отдельно, поскольку он влияет на форму спектральных кривых.

После вычисления преобразования Фурье обычно применяется операция \texttt{fftshift}: нулевая частота переносится в центр спектра. Тогда малые радиусы вокруг центра соответствуют низким частотам, а удаление от центра — переходу к более высоким частотам. Такое представление удобно для радиального анализа.

\subsection{Частотный радиус, круг Найквиста и радиальное представление}

Двумерный спектр содержит информацию и о масштабе частоты, и о её направлении. Например, горизонтальная и вертикальная граница могут иметь похожий частотный масштаб, но разные направления в частотной плоскости. В данной работе направление отдельно не анализируется. Меня интересует зависимость энергии от частотного радиуса, то есть от расстояния частотной точки до центра спектра.

Радиус частотной точки после \texttt{fftshift} определяется как расстояние до центра спектра:
\[
r(u,v)=\sqrt{(u-u_0)^2+(v-v_0)^2},
\]
где $(u_0,v_0)$ — координаты центра. Частоты с близкими значениями радиуса объединяются в кольца. Так двумерная карта спектра превращается в одномерный радиальный профиль.

Для карт признаков конечного размера максимальная представимая частота ограничена частотой Найквиста. В одномерном случае это половина частоты дискретизации. В двумерном случае ситуация чуть менее очевидна, потому что частотная плоскость имеет горизонтальное, вертикальное и диагональное направления. В работе используется радиус
\[
R_{\max}=\frac{\min(H,W)}{2}.
\]
Такой выбор соответствует анализу внутри круга Найквиста, вписанного в частотную плоскость. Это делает сравнение разных направлений более аккуратным: в расчёт не включаются диагональные области, которые выходят за пределы горизонтальной или вертикальной границы Найквиста.

Нормированный радиус задаётся как
\[
\rho=\frac{r}{R_{\max}}.
\]
После такой нормализации $\rho=0$ соответствует центру спектра, а $\rho=1$ — границе выбранного круга Найквиста. В результатах дополнительно используется масштабирование частотной оси к разрешению начального слоя, чтобы сравнивать слои с разным пространственным размером на общей оси. Собственно, без этого поздние слои ResNet и Swin нельзя было бы напрямую сопоставлять с ранними: у них просто другой доступный диапазон частот.

\subsection{Радиальное среднее и радиальная энергия}

В радиальном анализе есть два близких, но разных объекта. Первый — средняя мощность в кольце. Она получается, если суммарную мощность на радиусе разделить на число точек в этом кольце. Такая кривая хорошо подходит для визуализации: дальние кольца не становятся больше только из-за того, что в них больше частотных точек.

Второй объект — суммарная энергия кольца. Она нужна для метрик. Если требуется посчитать долю энергии в низких частотах, нельзя брать только среднее значение в кольце: кольца разного радиуса содержат разное число точек. Поэтому в методике для графиков используется радиальная средняя мощность, а для центроида и доли низкочастотной энергии — суммарная энергия кольца с последующей нормализацией. Это различие важно, потому что иначе величина, названная долей энергии, на самом деле не будет учитывать количество частотных компонент в диапазоне.

\subsection{Интерпретация низких и высоких частот}

В обычной обработке изображений низкие частоты связывают с крупной формой и плавными областями, а высокие — с мелкими деталями, границами, текстурами и шумом. Для внутренних представлений нейронных сетей эта интерпретация сохраняется только частично. Карта признаков — это уже не исходное изображение, а отклик некоторого набора фильтров или токенных преобразований. Поэтому низкая частота в карте признаков означает пространственно плавное распределение признака, а высокая — быстрое изменение этого признака по координатам.

По этой причине выводы должны быть осторожными. Если спектр глубокого слоя сдвигается к низким частотам, это может соответствовать более гладкому и агрегированному представлению. Но это не доказывает автоматически, что модель «лучше понимает объект» или что высокие частоты стали ненужными. Аналогично, широкий хвост спектра может отражать сохранение локальных различий, но не доказывает, что эти различия полезны для классификации.

Частотные компоненты действительно связаны с поведением моделей компьютерного зрения. Например, работа \cite{yin2019fourier} рассматривает устойчивость моделей через призму частотной области, а в \cite{wang2020highfrequency} обсуждается роль высокочастотных компонент в обобщающей способности CNN. Для данной диссертации эти работы используются как контекст. В самой экспериментальной части не проверяется причинная роль конкретных частот для качества модели. Проверяется другое: как архитектуры меняют спектральный профиль внутренних представлений.

\section{Архитектуры компьютерного зрения}

\subsection{Свёрточные нейронные сети и ResNet}

Свёрточная нейронная сеть — это архитектура, рассчитанная на данные с пространственной структурой. В изображении соседние пиксели обычно связаны друг с другом, и CNN использует это свойство через локальные фильтры. Один и тот же фильтр применяется ко всем пространственным позициям, поэтому сеть может находить похожие структуры в разных частях изображения.

Для понимания CNN нужны несколько терминов. Канал — отдельная карта признаков. Ядро свёртки — небольшой набор весов, который применяется к локальному участку входа. Шаг свёртки задаёт, как фильтр перемещается по координатам. Если шаг больше единицы, пространственное разрешение уменьшается. Понижение разрешения также может выполняться через пулинг. Рецептивное поле — это область исходного изображения, которая влияет на значение в конкретной позиции карты признаков.

С точки зрения спектрального анализа CNN естественно рассматривать как иерархическую систему локальных преобразований. Ранние уровни работают с более плотной пространственной сеткой. Поздние уровни часто имеют меньшее разрешение и большее рецептивное поле. Поэтому при переходе к глубоким слоям можно ожидать смещение спектральной энергии к более низким частотам. Но здесь есть методический риск: часть этого сдвига может быть простой геометрией, а не изменением самих признаков. Когда разрешение карты уменьшается, доступный диапазон частот тоже уменьшается. Именно поэтому в результатах используются контрольные эксперименты на понижение разрешения.

Семейство ResNet основано на остаточных связях \cite{he2016deep}. Остаточная связь позволяет блоку обучать не полное преобразование входа, а добавку к нему. За счёт этого глубокие сети легче оптимизируются, а градиент проходит через модель устойчивее. Для данной работы ResNet удобна как представитель классической иерархической CNN: внутри одного семейства доступны модели разной глубины — ResNet18, ResNet34, ResNet50 и ResNet101.

В экспериментальной части ResNet анализируется по четырём крупным стадиям: \texttt{layer1}, \texttt{layer2}, \texttt{layer3} и \texttt{layer4}. Это не отдельные операции, а группы блоков. Между стадиями меняются разрешение, число каналов и эффективный масштаб признаков. Поэтому результаты по ResNet нужно читать вместе с контролем: отдельно смотрится общая динамика по стадиям и отдельно проверяется, что происходит при искусственном уменьшении разрешения раннего слоя.

\subsection{Transformer и механизм самовнимания}

Transformer возник как архитектура для обработки последовательностей, в которой основной обмен информацией между элементами строится через механизм самовнимания \cite{vaswani2017attention}. В свёрточной сети локальность задаётся фильтром. В Transformer токены смешиваются через веса внимания, которые вычисляются из самих представлений.

Механизм самовнимания можно описать как адаптивное смешивание токенов. Для каждого токена модель оценивает, какие другие токены сильнее участвуют в его обновлении. Многоголовое внимание делает это в нескольких параллельных подпространствах. Но обычный блок трансформера не состоит только из внимания. В нём есть нормализация, MLP-блок и остаточные связи. Поэтому спектральное поведение трансформера нельзя выводить только из матрицы внимания.

В литературе самовнимание иногда рассматривают как оператор с низкочастотными свойствами. Например, в работе \cite{park2022how} показано, что внимание может вести себя как низкочастотный фильтр. Но переносить этот вывод на весь Vision Transformer (ViT)  напрямую нельзя. Полный блок содержит несколько операций, и их совместное действие может дать немонотонный спектральный профиль. В данной работе это как раз проверяется экспериментально: если у ViT спектр не сжимается монотонно, это не противоречит идее о сглаживающих свойствах внимания, потому что анализируется выход всего блока, а не изолированный оператор внимания.

\subsection{Vision Transformer}

Vision Transformer переносит трансформерный подход на изображения \cite{dosovitskiy2021image}. Изображение разбивается на непересекающиеся патчи. Например, при размере входа $224\times224$ и размере патча $16\times16$ получается решётка $14\times14$, то есть 196 патчей. Каждый патч разворачивается в вектор и проходит через линейную проекцию. После этого он становится токеном.

Чтобы модель различала расположение патчей, к токенам добавляются позиционные эмбеддинги. В классификационных версиях ViT обычно используется CLS-токен. Он не соответствует отдельному участку изображения и не является позиционным эмбеддингом; его задача — агрегировать информацию для классификационного выхода. В методике CLS-токен исключается перед построением спектра, потому что у него нет координаты в двумерной решётке патчей.

ViT отличается от CNN не только наличием внимания. У него другая начальная дискретизация изображения: один токен соответствует целому патчу. При этом внутри стандартного ViT размер патчевой решётки сохраняется одинаковым по блокам. Поэтому изменения спектра между блоками нельзя объяснить межстадийным уменьшением разрешения, как в ResNet или Swin. Это делает ViT удобным контрастным случаем для проверки: меняется ли спектральный профиль при фиксированной решётке.

Сравнение ViT и CNN подробно обсуждается в \cite{raghu2021do}. Авторы показывают, что визуальные трансформеры иначе распределяют информацию по глубине и раньше используют глобальные связи. Для данной работы это означает, что спектральная динамика ViT не обязана повторять поведение ResNet. Но в тексте результатов это формулируется как наблюдение по спектрам и метрикам, а не как доказательство конкретного механизма внимания.

\subsection{Swin Transformer}

Swin Transformer был предложен как иерархическая трансформерная архитектура для изображений \cite{liu2021swin}. В отличие от классического ViT, он не использует глобальное внимание по всем токенам на каждом блоке. Внимание вычисляется внутри локальных окон, а на соседних блоках окна сдвигаются. Такой механизм называется оконным смещением (shifted windows).

Ещё один элемент Swin — слияние патчей (patch merging). Эта операция объединяет соседние патчи, уменьшает пространственное разрешение и увеличивает число каналов. По смыслу она похожа на понижение разрешения в CNN. Благодаря этому Swin строит многостадийную иерархию: ранние стадии имеют более плотную решётку, поздние — меньшую решётку и более агрегированные признаки.

С точки зрения спектрального анализа Swin интересен именно как промежуточный случай. С одной стороны, он использует токены и внимание. С другой стороны, у него есть стадии и уменьшение разрешения, как у CNN. Поэтому можно ожидать, что его финальные представления будут более низкочастотными. Однако объяснять это только оконным вниманием или сдвигом окон нельзя. В данной работе не проводится анализ вклада, где отдельно отключаются оконные сдвиги или слияние патчей. Поэтому в результатах такие объяснения используются только как возможная интерпретация. Проверенный факт — форма спектральных кривых и значения метрик.

\section{Спектральные свойства нейронных сетей}

\subsection{Спектральное смещение и принцип частот}

Одно из направлений в теории глубокого обучения связано с понятием spectral bias, или спектрального смещения. Под этим понимается склонность нейронных сетей легче осваивать низкочастотные зависимости, чем высокочастотные. Иначе говоря, более гладкие компоненты функции часто усваиваются раньше, а более резкие зависимости оказываются сложнее.

Работа \cite{rahaman2019spectral} является одной из основных в этом направлении. В ней показано, что глубокие сети с ReLU-активациями обладают смещением в сторону низких частот. Это не означает, что сеть не может выучить высокочастотную функцию. Скорее речь о том, что низкочастотные закономерности оказываются для неё более естественными в процессе обучения.

Эта идея развивалась в рамках frequency principle \cite{xu2025overview}. В работе \cite{kiessling2022computable} рассматриваются способы вычислимого определения спектрального смещения, то есть переход от качественного описания к более формальному измерению. Работа \cite{tancik2020fourier} показывает, что Фурье-признаки могут помогать сетям лучше представлять высокочастотные функции.

Для данной диссертации эти работы задают общий фон. Но есть принципиальная разница: spectral bias обычно говорит о процессе обучения функции, а в данной работе анализируются уже обученные модели и их внутренние представления. Поэтому результаты не следует напрямую отождествлять с динамикой обучения. Здесь исследуется не то, какие частоты модель учит первыми, а то, как распределена энергия признаков в готовых представлениях.

\subsection{Частоты, устойчивость и трансформерные модели}

Частотная область часто используется для анализа устойчивости и поведения моделей компьютерного зрения. В работе \cite{yin2019fourier} искажения рассматриваются в частотной области, что позволяет изучать чувствительность моделей к изменениям разных масштабов. В \cite{wang2020highfrequency} обсуждается роль высокочастотных компонент в обобщающей способности CNN. Эти работы показывают, что частотный состав входа и признаков может быть связан с поведением модели.

Для визуальных трансформеров вопрос частот также обсуждается в нескольких работах. Исследования \cite{naseer2021intriguing,paul2022vision,alijani2024vision} показывают, что ViT могут отличаться от CNN по внутренней организации признаков и устойчивости. Работа \cite{bai2022improving} отдельно рассматривает высокочастотные компоненты в ViT. Всё это поддерживает идею, что CNN и трансформеры стоит сравнивать не только по итоговой точности, но и по структуре внутренних представлений.

При этом в данной работе не проверяется устойчивость и не проводится оценка причинной роли отдельных частот. Частотные исследования используются как контекст, а экспериментальная часть ограничивается спектральным анализом промежуточных представлений.

\subsection{Частотно-ориентированные методы в глубоком обучении}

Частотная область используется не только для анализа, но и для построения новых методов. В работе \cite{xu2020learning} рассматривается обучение в частотной области. Работа \cite{li2021frequency} показывает использование частотного домена в CNN для задач диагностики и адаптации домена. В более широком контексте глубокое обучение активно применяется к спектральному анализу, моделированию и восстановлению сигналов \cite{liu2024deep}.

Для трансформеров также появляются частотно-ориентированные архитектуры и модификации. Обзоры \cite{aslan2026frequency,li2025review} рассматривают модели, использующие преобразования Фурье, вейвлеты, дискретное косинусное преобразование и другие инструменты частотной области в задачах обработки изображений.

Для данной работы это важно как методический контекст. Частотная область уже используется в глубоких моделях, но здесь она применяется не для изменения архитектуры, а для диагностики внутренних представлений. Иными словами, работа отвечает не на вопрос «как улучшить модель», а на более узкий вопрос: как разные архитектуры перераспределяют спектральную энергию по глубине.

\subsection{Набор данных ImageNet и ImageNet 10K}

ImageNet является одним из основных наборов изображений в истории компьютерного зрения \cite{deng2009imagenet}. Он содержит большое число изображений и классов, а многие современные архитектуры обучаются или оцениваются именно на ImageNet. Поэтому этот набор данных удобен для сопоставления предобученных моделей.

В данной работе используется подвыборка ImageNet 10K \cite{priyerana_imagenet10k}. Она содержит по 10 изображений для каждого из 1000 классов ImageNet. Такой набор сохраняет разнообразие классов, но остаётся вычислительно выполнимым для спектрального анализа многих моделей и промежуточных блоков.

Здесь ImageNet 10K используется не для обучения и не для проверки итоговой точности. Его роль другая: подать на предобученные модели одинаковые естественные изображения и сравнить их внутренние представления. Поэтому далее в работе точность (accuracy) не рассматривается как отдельная зависимая переменная. Основной предмет исследования — спектральная организация признаков.

\section{Выводы по теоретической главе}

В этой главе введены понятия, которые нужны для методики и результатов: внутреннее представление, карта признаков, токен, патч, спектр мощности, DC-компонента, частотный радиус, круг Найквиста, радиальное среднее и радиальная энергия. Это нужно для того, чтобы дальше не смешивать разные вещи: визуальную радиальную кривую и численную долю энергии, карту признаков CNN и решётку токенов ViT, реальное изменение признаков и эффект уменьшения разрешения.

ResNet, ViT и Swin рассматриваются как три разных способа обработки пространственного сигнала. ResNet представляет свёрточную иерархию с понижением разрешения между стадиями. ViT работает с фиксированной решёткой патчевых токенов. Swin совмещает токенную обработку с иерархией и слиянием патчей. Из-за этих различий можно ожидать разные спектральные траектории по глубине.

Дальнейшая методика строится именно на этом различии. Для ResNet и Swin требуется контроль влияния понижения разрешения, потому что уменьшение разрешения само по себе сужает доступный частотный диапазон. Для ViT такой межстадийный контроль не нужен в том же виде, поскольку размер патчевой решётки сохраняется по блокам. В результатах наблюдения по графикам отделяются от возможных интерпретаций через внимание, оконную обработку или слияние патчей, так как причинные механизмы отдельно не отключались и не проверялись.
